% ch02 Non-interacting electron gas (书页 10–15 / 分册p024–p029)
% 第2章 Non-interacting electron gas
\chapter{无相互作用电子气}

在上一章的结尾，我们指出，可以用无相互作用的电子与声子来理解金属的基本性质。例如，金属的低温比热容是两项之和：来自电子的、随温度 $T$ 线性增长的项，以及来自声子的、正比于 $T^{3}$ 的项。这一结果可以从无相互作用粒子图像得出。由非磁性杂质散射所限制的电导率，同样可以很好地用无相互作用电子气来描述。此外，凭借单电子能带理论的知识，人们便可以定性地辨别金属、绝缘体与半导体之间的差别。无相互作用模型取得的显著成功是自相矛盾的，因为电子和离子既与自身强烈相互作用，彼此之间也强烈相互作用。在成功之余，无相互作用图像也存在着巨大的缺陷，最明显的是它无法描述一些老问题，如内聚能、超导电性、磁性，也无法描述一些较新的现象，如掺杂莫特（Mott）绝缘体、近藤（Kondo）问题以及分数量子霍尔效应。我们首先回顾无相互作用电子气的物理。只有在我们发展出处理电子相互作用的方法之后，才可能阐明无相互作用模型为何如此有效的原因。

金属中的电子是自旋为 $\hbar/2$、服从费米–狄拉克（Fermi--Dirac）统计的量子力学粒子。单个电子的哈密顿量为 $\widehat{p}^{2}/2m$，其中 $\widehat{\mathbf{p}}$ 是电子动量（算符），$m$ 是电子质量。它的本征态是形式为 $\mathrm{e}^{\mathrm{i}\mathbf{p}\cdot r/\hbar}/\sqrt{V}$ 的平面波乘以一个旋量，该旋量给出电子自旋在某个方便选取的轴（通常取 $\hat{z}$）上的投影 $\hbar\sigma/2$，其中 $\sigma=\pm1$；这里 $V$ 是系统体积。$N$ 个这种无相互作用电子的哈密顿量（算符）
\begin{equation*}
\widehat{H}=\sum_{i=1}^{N}\frac{\widehat{p}_{i}^{2}}{2m},
\tag{2.1}
\end{equation*}
就是各个粒子动能之和。此时，本征态是被占据的那些单粒子平面波态的乘积。每个平面波态至多可被一个给定自旋的电子占据。我们用分布函数 $f_{\mathbf{p}\sigma}$ 来标记这些本征态：若单粒子动量–自旋态被占据，则 $f_{\mathbf{p}\sigma}$ 取 1，否则取 0。在基态中，最低的 $N/2$ 个单粒子态被自旋相反的电子双重占据。因此，在基态（温度 $T=0$）下，分布函数为
\begin{equation*}
f_{\mathbf{p}\sigma}=\Theta(\mu_{0}-p^{2}/2m),
\tag{2.2}
\end{equation*}
其中 $\Theta(x)$ 是亥维赛（Heaviside）函数，当 $x>0$ 时 $\Theta(x)=1$，否则为 0。这里 $\mu_{0}$ 是零温电子化学势，在此情形下它就是费米能，即最高被占据态的能量 $p_{\mathrm{F}}^{2}/2m$，其中 $p_{\mathrm{F}}$ 是电子费米动量。费米温度 $T_{\mathrm{F}}$ 等于 $\mu_{0}/k_{\mathrm{B}}$。

用 $f_{\mathbf{p}\sigma}$ 表示，电子的总数为
\begin{equation*}
N=\sum_{\mathbf{p},\sigma}f_{\mathbf{p}\sigma}.
\tag{2.3}
\end{equation*}
在基态中，我们可以把求和换成积分，求得 $T=0$ 时的电子密度：
\begin{equation*}
n_{\mathrm{e}}=\frac{N(T=0)}{V}=\frac{2}{V}\sum_{p<p_{\mathrm{F}}}=2\int_{0}^{p_{\mathrm{F}}}\frac{\mathrm{d}\mathbf{p}}{(2\pi\hbar)^{3}}=\frac{p_{\mathrm{F}}^{3}}{3\pi^{2}\hbar^{3}}.
\tag{2.4}
\end{equation*}
粒子间的平均间距基本上就是容纳单个电子的球体的半径 $r_{\mathrm{e}}$，即
\begin{equation*}
\frac{4\pi r_{\mathrm{e}}^{3}}{3}n_{\mathrm{e}}=1.
\tag{2.5}
\end{equation*}
于是，由式 (2.4)，粒子间距的尺度为
\begin{equation*}
r_{\mathrm{e}}=\left(\frac{9\pi}{4}\right)^{1/3}\frac{\hbar}{p_{\mathrm{F}}}=1.92\,\frac{\hbar}{p_{\mathrm{F}}},
\tag{2.6}
\end{equation*}
它与晶格间距同数量级。当 $r_{\mathrm{e}}\approx 1\,\text{\AA}$ 时，我们得到费米速度 $v_{\mathrm{F}}=p_{\mathrm{F}}/m\approx\hbar/mr_{\mathrm{e}}\approx 10^{8}\,\mathrm{cm/s}\approx c/300$，其中 $c$ 是光速。（对于基态中电子的运动，相对论效应一般并不重要。）惯常的做法是采用无量纲比值 $r_{\mathrm{s}}=r_{\mathrm{e}}/a_{0}$，其中 $a_{0}=\hbar^{2}/me^{2}$ 是玻尔（Bohr）半径；这个量提供了电子密度的一种量度。稠密极限对应于 $r_{\mathrm{s}}\ll 1$，稀薄区则对应于 $r_{\mathrm{s}}\gg 1$。金属的 $r_{\mathrm{s}}$ 在 2 到 6 之间变化。表 2.1 列出了碱金属的 $r_{\mathrm{s}}$ 值。事实上，铯是所有金属中最稀薄的。正是这个大的 $r_{\mathrm{s}}$ 值导致了 Cs 密度的不均匀性。在第 5 章中，我们将进一步讨论与大 $r_{\mathrm{s}}$ 相关的物理，例如最终会形成维格纳（Wigner）晶体。我们还可以利用式 (2.4) 求解零温化学势
\begin{equation*}
\mu_{0}=(3\pi^{2})^{2/3}\frac{\hbar^{2}}{2m}n_{\mathrm{e}}^{2/3}.
\tag{2.7}
\end{equation*}
对 Na，$\mu_{0}=3.1\,\mathrm{eV}$；典型金属的 $\mu_{0}$ 介于 1 与 5 eV 之间。

\begin{table}[H]
\centering
\caption*{表 2.1\quad 碱金属的 $r_{\mathrm{s}}$ 值}
\begin{tabular}{lccccc}
\toprule
 & Li & Na & K & Rb & Cs\\
\midrule
$r_{\mathrm{s}}$ & 3.25 & 3.93 & 4.86 & 5.2 & 5.62\\
\bottomrule
\end{tabular}
\end{table}

系统的总能量是对被占据态求和、并以单粒子能量 $\epsilon_{\mathbf{p}}=p^{2}/2m$ 加权的结果：
\begin{equation*}
E=\sum_{\mathbf{p},\sigma}\epsilon_{\mathbf{p}}f_{\mathbf{p}\sigma}.
\tag{2.8}
\end{equation*}
在 $T=0$ 时，能量为
\begin{equation*}
E_{0}=\frac{p_{\mathrm{F}}^{5}V}{10m\pi^{2}\hbar^{3}}=\frac{3}{5}N\mu_{0}.
\tag{2.9}
\end{equation*}
由热力学关系 $P=-(\partial E/\partial V)_{T,N}$，我们求得基态的压强 $P_{0}=2\mu_{0}n_{\mathrm{e}}/5$。这个量是 $10^{6}$ atm 的量级，完全来源于粒子之间的不相容原理。

在有限温度下，我们定义一个在 0 与 1 之间取值、量度单粒子态平均占据情况的分布函数。对于在化学势 $\mu$ 下处于平衡的系统，
\begin{equation*}
f_{\mathbf{p}\sigma}=\frac{1}{\mathrm{e}^{\beta(\epsilon_{\mathbf{p}}-\mu)}+1}
\tag{2.10}
\end{equation*}
就是费米–狄拉克分布，其中 $\beta=1/k_{\mathrm{B}}T$。费米–狄拉克分布函数使给定能量和电子数下的熵取极大。一般地说，熵 $S$ 由与系统宏观热力学态相一致的微观状态数 $W$ 的对数给出：
\begin{equation*}
S=k_{\mathrm{B}}\ln W,
\tag{2.11}
\end{equation*}
对于电子问题，微观态由动量 $\mathbf{p}$ 标记，占据数为 0 或 1。然而，分布函数 $f_{\mathbf{p}\sigma}$ 是定义在全部动量态上的光滑函数。为构造这个函数，我们把动量态分成若干相格（cell），每个相格包含 $g_{i}$ 个动量态和 $n_{i}$ 个粒子。因为每个相格含有 $g_{i}$ 个态，
\begin{equation*}
\sum_{i}g_{i}\cdots=\sum_{\mathbf{p},\sigma}\cdots\,.
\tag{2.12}
\end{equation*}
对第 $i$ 个相格，把 $n_{i}$ 个粒子分布到 $g_{i}$ 个态中的不同方式的数目由组合因子 $W_{i}=g_{i}!/n_{i}!(g_{i}-n_{i})!$ 给出。对 $W_{i}$ 应用斯特林（Stirling）近似
\begin{equation*}
\ln N!\approx N(\ln N-1)
\tag{2.13}
\end{equation*}
后，我们得到
\begin{align*}
\ln W_{i}&\approx -n_{i}\ln\frac{n_{i}}{g_{i}}-(g_{i}-n_{i})\ln\frac{g_{i}-n_{i}}{g_{i}}
\tag{2.14}\\
&=-g_{i}\left[\frac{n_{i}}{g_{i}}\ln\frac{n_{i}}{g_{i}}+\left(1-\frac{n_{i}}{g_{i}}\right)\ln\left(1-\frac{n_{i}}{g_{i}}\right)\right],
\tag{2.15}
\end{align*}
其中 $n_{i}/g_{i}$ 是第 $i$ 个相格中态被占据的比例。实际上，$n_{i}/g_{i}=f_{i}$ 正是我们要找的光滑分布函数。若把这个表达式代入 $\ln W_{i}$ 的方程，便得到熵的熟知结果：
\begin{equation*}
S=-k_{\mathrm{B}}\sum_{\mathbf{p},\sigma}\left[f_{\mathbf{p}\sigma}\ln f_{\mathbf{p}\sigma}+(1-f_{\mathbf{p}\sigma})\ln(1-f_{\mathbf{p}\sigma})\right],
\tag{2.16}
\end{equation*}
其中我们已经利用式 (2.12) 把对相格的求和换成了对自旋和动量态的求和。为求得分布函数，我们在粒子数和能量固定的约束下使熵极大。将 $S-E/T+\mu N/T$ 对 $f_{\mathbf{p}\sigma}$ 取极值，我们得到
\begin{equation*}
0=\epsilon_{\mathbf{p}\sigma}-\mu-k_{\mathrm{B}}T\ln\left(f_{\mathbf{p}\sigma}^{-1}-1\right),
\tag{2.17}
\end{equation*}
这表明 $f_{\mathbf{p}\sigma}$ 就是式 (2.10) 中的费米分布函数。

现在让我们从热力学关系
\begin{equation*}
C_{V}=T\left(\frac{\partial S}{\partial T}\right)_{V}
\tag{2.18}
\end{equation*}
来计算热容。为计算熵对温度的导数，我们考虑熵的一般变分
\begin{align*}
\delta S&=-2k_{\mathrm{B}}\sum_{\mathbf{p}}\delta f_{\mathbf{p}}\ln\frac{f_{\mathbf{p}}}{1-f_{\mathbf{p}}}=2k_{\mathrm{B}}\sum_{\mathbf{p}}\delta f_{\mathbf{p}}\frac{\epsilon_{\mathbf{p}}-\mu}{k_{\mathrm{B}}T}
\tag{2.19}\\
&=2Vk_{\mathrm{B}}\int\frac{\mathrm{d}\mathbf{p}}{(2\pi\hbar)^{3}}\,\delta f_{\mathbf{p}}\,\frac{\epsilon_{\mathbf{p}}-\mu}{k_{\mathrm{B}}T},
\tag{2.20}
\end{align*}
即分布函数 $f_{\mathbf{p}\sigma}$ 的变分。通过引入单位体积的单粒子态密度，我们可以把这个表达式进一步化简：
\begin{align*}
N(\epsilon)&=2\int\frac{\mathrm{d}\mathbf{p}}{(2\pi\hbar)^{3}}\,\delta(\epsilon-\epsilon_{\mathbf{p}})
\tag{2.21}\\
&=\frac{1}{\pi^{2}\hbar^{3}}\int p^{2}\frac{\mathrm{d}p}{\mathrm{d}\epsilon_{\mathbf{p}}}\,\delta(\epsilon-\epsilon_{\mathbf{p}})\,\mathrm{d}\epsilon_{\mathbf{p}}
\tag{2.22}
\end{align*}
\begin{equation*}
N(\epsilon_{\mathbf{p}})=\frac{mp}{\pi^{2}\hbar^{3}}.
\tag{2.23}
\end{equation*}
这里的关键在于，单粒子态密度是动量的线性函数；这个量也可写为 $N(\epsilon_{\mathbf{p}})=(\mathrm{d}p/\mathrm{d}\epsilon_{\mathbf{p}})(p^{2}/\pi^{2}\hbar^{3})$。如果我们用态密度改写 $\delta S$，就得到
\begin{equation*}
\delta S=\frac{V}{T}\int\mathrm{d}\epsilon_{\mathbf{p}}\,\delta f_{\mathbf{p}}\,N(\epsilon_{\mathbf{p}})(\epsilon_{\mathbf{p}}-\mu),
\tag{2.24}
\end{equation*}
一个可以用索末菲（Sommerfeld）展开来计算的单一积分。尽管这个展开是标准方法，我们仍将回顾一下。

\section*{索末菲（Sommerfeld）展开}
\markright{索末菲（Sommerfeld）展开}
\addcontentsline{toc}{section}{索末菲（Sommerfeld）展开}

考虑如下形式的积分
\begin{equation*}
I=\int_{0}^{\infty}\mathrm{d}\epsilon\,f(\epsilon)h(\epsilon),
\tag{2.25}
\end{equation*}
其中 $h(\epsilon)$ 是任意光滑函数，而 $f(\epsilon)=1/(\mathrm{e}^{\beta(\epsilon-\mu)}+1)$。我们对 $I$ 分部积分：
\begin{equation*}
I=\int_{0}^{\infty}f'(\epsilon)H(\epsilon)\,\mathrm{d}\epsilon,
\tag{2.26}
\end{equation*}
其中
\begin{equation*}
H=-\int_{0}^{\epsilon}h(x)\,\mathrm{d}x.
\tag{2.27}
\end{equation*}
因为 $f'(\epsilon)$ 在化学势处有很陡的峰，我们可以在 $\epsilon=\mu$ 附近把 $H$ 作泰勒展开：
\begin{equation*}
H(\mu)+(\epsilon-\mu)\left(\frac{\partial H}{\partial\epsilon}\right)_{\epsilon=\mu}+\frac{1}{2}(\epsilon-\mu)^{2}\left(\frac{\partial^{2}H}{\partial\epsilon^{2}}\right)_{\epsilon=\mu}+\cdots\,,
\tag{2.28}
\end{equation*}
于是得到一列如下形式的、需要计算的积分
\begin{equation*}
L_{j}=-\int_{0}^{\infty}(\epsilon-\mu)^{j}f'(\epsilon)\,\mathrm{d}\epsilon
\tag{2.29}
\end{equation*}
对 $j=0$，$L_{0}=-f(\infty)+f(0)=1$。对其余的积分，我们可以把下限换成 $-\infty$。令 $x=\beta(\epsilon-\mu)$，我们有
\begin{equation*}
L_{j}=\frac{1}{\beta^{j}}\int_{-\infty}^{\infty}x^{j}\frac{\mathrm{e}^{x}}{(\mathrm{e}^{x}+1)^{2}}\,\mathrm{d}x.
\tag{2.30}
\end{equation*}
由于 $j$ 为奇数时被积函数是 $x$ 的奇函数，只有偶数 $j$ 的积分不为零。$L_{j}$ 的头几个值为
\begin{gather*}
L_{0}=1,
\tag{2.31}\\
L_{2}=\frac{\pi^{2}}{3}(k_{\mathrm{B}}T)^{2},
\tag{2.32}\\
L_{4}=\frac{7\pi^{4}}{15}(k_{\mathrm{B}}T)^{4}.
\tag{2.33}
\end{gather*}
于是，我们系统地得到 $I$ 的级数展开：
\begin{equation*}
I=\int_{0}^{\mu}h(\epsilon)\,\mathrm{d}\epsilon+\frac{\pi^{2}}{6}(k_{\mathrm{B}}T)^{2}h'(\mu)+\frac{7\pi^{4}}{360}(k_{\mathrm{B}}T)^{4}h''(\mu)+\cdots\,.
\tag{2.34}
\end{equation*}
对不依赖于温度的函数 $h$，展开的第一项与温度无关，因此在固定 $\mu$ 下，
\begin{equation*}
\int_{0}^{\infty}\mathrm{d}\epsilon\,\delta f(\epsilon)h(\epsilon)=\frac{\pi^{2}}{3}h'(\epsilon=\mu)k_{\mathrm{B}}^{2}T\delta T
\tag{2.35}
\end{equation*}
是 $f_{\mathbf{p}}$ 随温度变化时的领头项。

现在让我们回到低温熵的计算。我们要计算固定粒子数下熵随温度的变化。我们先在固定 $\mu$ 下计算，然后再证明：固定粒子数时化学势随温度的变化在这里可以忽略。如果把式 (2.35) 代入式 (2.24)，我们发现单位体积的熵的变分
\begin{equation*}
\delta s=\frac{\pi^{2}}{3}k_{\mathrm{B}}^{2}N(\epsilon_{\mathrm{F}})\delta T,
\tag{2.36}
\end{equation*}
是一个不依赖于温度的常数。单位体积的热容
\begin{equation*}
c_{V}=T\left(\frac{\delta s}{\delta T}\right)_{V}=\frac{\pi^{2}}{3}k_{\mathrm{B}}^{2}N(\epsilon_{\mathrm{F}})T,
\tag{2.37}
\end{equation*}
按温度的线性函数标度。这一贡献完全来自传导电子。每个电子的贡献为
\begin{equation*}
\frac{c_{V}}{n_{\mathrm{e}}}=\frac{\pi^{2}k_{\mathrm{B}}^{2}mT}{p_{\mathrm{F}}^{2}}=\frac{\pi^{2}}{2}k_{\mathrm{B}}\frac{T}{T_{\mathrm{F}}}.
\tag{2.38}
\end{equation*}
把这个结果与经典热容 $3k_{\mathrm{B}}/2$ 比较，我们发现量子力学的数值要小一个因子 $\pi^{2}T/3T_{\mathrm{F}}$。在金属中，只有一部分电子处在费米能级上。比值 $T/T_{\mathrm{F}}$ 给出的正是能量处于费米能附近 $k_{\mathrm{B}}T$ 范围内的那部分电子的比例。电子在费米能级以下越深，它对热容的贡献就越小。

我们可以用索末菲展开证明，作为密度和温度的函数，化学势在低温 $T$ 下由下式给出：
\begin{equation*}
\mu(n_{\mathrm{e}},T)=\mu(n_{\mathrm{e}},0)\left(1-\frac{\pi^{2}}{12}\left(\frac{T}{T_{\mathrm{F}}}\right)^{2}\right),
\tag{2.39}
\end{equation*}
其中 $\mu(n_{\mathrm{e}},0)=\epsilon_{\mathrm{F}}$。这个结果的证明与推导熵时所用的方法类似，我们把它留作习题（习题 2.2）。由式 (2.39) 可见，固定 $n_{\mathrm{e}}$ 时化学势的第一修正是 $T^{2}$ 量级的，因此在计算低温熵，即式 (2.36) 时，这一修正可以忽略。

\section*{习题}
\addcontentsline{toc}{section}{习题}

\textbf{习题 2.1}\quad 计算积分
\begin{equation*}
L_{j}=\int_{-\infty}^{\infty}x^{j}\frac{\mathrm{e}^{x}}{(\mathrm{e}^{x}+1)^{2}}\,\mathrm{d}x.
\end{equation*}
并特别证明，只有偶数 $j$ 的积分才不为零。

\textbf{习题 2.2}\quad 利用索末菲展开计算费米气体化学势的第一温度修正（式 (2.39)），把它表示为密度的函数。
